\documentclass[11pt]{article}

\usepackage[margin=1in]{geometry}
\usepackage{amsmath,amssymb}
\usepackage{booktabs}
\usepackage{xcolor}
\usepackage{hyperref}

\title{A Rights-Safe Local Soccer Analytics Pipeline:\\From Licensed Footage to Match-Outcome Prediction}

\author{
  The Soccer-Edge Collaboration\\
  \texttt{github.com/scottjoyner/soccer-analytics}
}
\date{Draft --- July 2026}

\begin{document}

\maketitle

\begin{abstract}
We present an open, reproducible pipeline for soccer analytics that enforces
data-rights compliance by construction while turning permitted \emph{local}
footage and open event data into structured tables, detection datasets, and
match-outcome models. The system ingests open event feeds (StatsBomb Open Data,
Metrica Sports sample data, and approved SoccerNet subsets) together with
rights-approved local video, runs a YOLO-based computer-vision abstraction to
detect players and the ball, aggregates per-player and per-team statistics, and
fine-tunes match-outcome predictors (winner classification and home/away score
regression) as well as a convolutional occupancy-grid model. A lineage-and-rights
gate refuses to process any footage that lacks a recorded written-rights
reference, so no scraped or unlicensed audiovisual content can enter the
pipeline. We report a reproducible case study over 98 permitted FIFA World Cup
2026 highlight clips, yielding 88,127 detection rows, and we discuss the
honest limitation that short highlight-clip detection aggregates carry little
predictive signal for match outcome.
\end{abstract}

\section{Introduction}
Data-driven soccer analytics has advanced rapidly, but most pipelines assume
unfettered access to broadcast footage and event feeds. For academic research
this assumption is increasingly untenable: audiovisual rights are tightly held,
and scraping or caching copyrighted video violates both platform terms and the
data-governance expectations of institutional review. We argue that a
\emph{rights-safe, local-first} design is not merely a compliance footnote but a
first-class system property that improves reproducibility: if the pipeline can
only consume data the researcher is entitled to process, then every published
result is, by construction, re-runnable by others with equivalent rights.

This paper describes \textsc{Soccer-Edge}, a pipeline organized around three
invariants: (i) public video URLs are \emph{discovery metadata only} and are
never processing inputs; (ii) every processed table and derived feature table
carries lineage fields and an explicit \texttt{rights\_reference}; and (iii) a
single rights gate is enforced at every footage entry point. We validate the
system end-to-end on 98 permitted local highlight clips and open event data.

\section{Related Work}
Event-stream datasets such as StatsBomb Open Data~\cite{statsbomb}, Metrica
Sports sample tracking~\cite{metrica}, and SoccerNet~\cite{soccernet} have
enabled reproducible spatiotemporal modeling without audiovisual rights. Prior
computer-vision work detects players, the ball, and team possession from
broadcast video using object detectors such as YOLO~\cite{redmon2016yolo,
ultralytics} and tracks them across frames. Our contribution is not a new
detector but a \emph{governance wrapper}: a reproducible, rights-gated pipeline
that composes existing open components (Ultralytics YOLO, scikit-learn,
PyTorch) into a research workflow that provably cannot ingest unlicensed media.

\section{System Design and the Rights Model}
The pipeline models provenance explicitly. Every processed table records
\texttt{lineage\_source\_name}, \texttt{lineage\_source\_path},
\texttt{lineage\_dataset\_version}, and \texttt{lineage\_ingested\_at\_utc}.
Footage rows carry a \texttt{rights\_status} drawn from
$\{\texttt{owned}, \texttt{licensed}, \texttt{compatible\_license},
\texttt{pending}, \texttt{blocked}\}$, and any approved status \emph{requires} a
non-empty \texttt{rights\_reference} (an explicit written-rights pointer such as
a permit file or a recorded consent). 

A single assertion, \texttt{assert\_processable}, is reused across the
command surface: \texttt{discover video} (metadata only), \texttt{catalog-local},
\texttt{video plan}, \texttt{detect-yolo}, \texttt{process},
\texttt{process-local-model}, \texttt{export-frames},
\texttt{train player-ball}, and \texttt{train local-finetune}. When a
manifest row is named, the footage file is opened only if its
\texttt{rights\_reference} is recorded \emph{and} its \texttt{local\_path}
matches the input. This makes the rights check a hard, test-covered gate rather
than documentation.

\section{Data}
\paragraph{Open event data.} StatsBomb Open Data, Metrica sample events/tracking,
and approved SoccerNet subsets are ingested into normalized Parquet tables with
lineage. From these we derive per-match player statistics (shots, passes,
carries, pressures, interceptions, tackles, fouls) and cross-match aggregates
with per-player totals, per-match averages, and team/opponent splits.

\paragraph{Permitted local footage.} For this study we use 98 FIFA World Cup 2026
highlight clips that were provided to the authors as permitted local files.
Their match results (teams, scores, penalty shootouts) are encoded directly in
the filenames and were parsed into a structured results table; no audiovisual
content was scraped or downloaded.

\paragraph{Codec handling.} The provided clips are AV1-encoded. The local media
reader has no AV1 decoder, so each clip is first transcoded to H.264 with
FFmpeg and then passed to the detector. This transcode step is a
rights-preserving local transformation of already-permitted footage; the
original files remain the system of record and the \texttt{rights\_reference}
propagates to every derived table.

\section{Computer-Vision Abstraction}
Detections are produced by a \texttt{YOLODetector} wrapping Ultralytics
YOLOv8n~\cite{ultralytics}, exposed through the \texttt{detect-yolo} command.
For each clip it emits per-frame player/ball bounding boxes, frame-to-frame
tracks, and pitch-space ball/player states. From these we compute, per match,
aggregate features
\[
f = (\texttt{n\_player},\ \texttt{n\_ball},\ \texttt{avg\_det\_per\_frame},\
\texttt{ball\_center\_x},\ \texttt{ball\_center\_y})
\]
normalized to the frame. We additionally rasterize each frame into an
occupancy grid (channels $\times$ height $\times$ width, default
$3\times 8\times 8$) for the convolutional model.

\section{Feature Aggregation}
Player statistics are aggregated with \texttt{build\_player\_aggregates}, which
yields cross-match totals, per-match averages, appearances, and an overall pass
completion rate, with optional \texttt{split\_by=team} and
\texttt{split\_by=opponent} breakdowns. For the highlight study, filename parsing
produced a 98-match results table (50 distinct teams) that supplies the
supervisory labels (winner, home/away score) for the outcome models.

\section{Match-Outcome Models}
\paragraph{Tabular model.} \texttt{train\_match\_predictor} fits a multinomial
logistic-regression winner classifier and two RandomForest regressors for the
home and away scores, using the five detection-derived features above. To avoid
optimistic in-sample estimates we evaluate behind a leakage-safe stratified
train/test split (68/30), a \texttt{StandardScaler}, and a
\texttt{CalibratedClassifierCV} wrapper so the winner probabilities are
calibrated and the Brier score is meaningful.
\paragraph{Convolutional model.} \texttt{train\_match\_predictor\_cnn} builds an
NPZ of occupancy-grid tensors (grouped by match, with a sliding window of
sequence length 4) and trains a small CNN winner classifier (PyTorch).

\section{Case Study and Results}
Over the 98 permitted clips the batch produced \textbf{88,127} detection rows
(86,007 player detections, 1,541 ball detections) across 7,878 grid frame-rows.
The tabular model was trained on the resulting 98-match dataset.

We compare two feature sets: (i) the five \emph{count} features used above
(\texttt{n\_player}, \texttt{n\_ball}, \texttt{avg\_det\_per\_frame},
\texttt{ball\_center\_x/y}), and (ii) nine \emph{track} features that attempt to
capture possession and pressure from the per-frame boxes---ball frame-rate,
person frame-rate, mean nearest player--ball distance, mean players-near-ball,
contested-ball rate, ball movement, and mean player-box area. Both are evaluated
behind the same stratified 68/30 split with a calibrated winner classifier.

\begin{table}[h]
\centering
\caption{Out-of-sample results on the 98-match highlight dataset (30-match held-out
test, calibrated classifier). Accuracy is reported against the majority-class
baseline (49\%); a Brier score near 0.5--0.7 indicates near-random probability
quality.}
\label{tab:results}
\begin{tabular}{lccc}
\toprule
Feature set / metric & Count & Track & Note \\
\midrule
Winner accuracy (test) & 53.3\% & 50.0\% & baseline 49.0\% \\
Winner Brier (test) & 0.623 & 0.614 & lower is better \\
Home-score MSE (test) & 1.770 & 1.864 & out-of-sample \\
Away-score MSE (test) & 2.345 & 2.066 & out-of-sample \\
Train rows / test rows & \multicolumn{2}{c}{68 / 30} & stratified \\
Detection rows & \multicolumn{2}{c}{88,127} & 86,007 player / 1,541 ball \\
Grid frame-rows & \multicolumn{2}{c}{7,878} & CNN input \\
\bottomrule
\end{tabular}
\end{table}

\paragraph{Honest limitation.} As Table~\ref{tab:results} shows, neither the
count nor the track detection features outperform the majority-class baseline
for winner prediction, and the calibrated Brier scores ($\approx0.61$) confirm
the probabilities are no better than chance. This is the expected and, we argue,
informative result: short highlight reels are highlight-selected (not
representative of open play), contain far fewer frames than a full match, and
aggregate detection counts are nearly invariant to final score. The pipeline
therefore demonstrates end-to-end reproducibility and governance---including a
leakage-safe, calibrated evaluation---while clearly delineating where predictive
signal must come from: full-match footage or rich open event data
(Section~\ref{sec:data}). We report out-of-sample figures only and do not claim
generalization.

\section{Reproducibility}
The study is scripted and re-runnable: \texttt{scripts/process\_highlights\_batch.py}
(transcode + detect, resumable), \texttt{scripts/build\_highlights\_training.py}
(aggregate + fine-tune), and \texttt{scripts/evaluate\_highlights.py} (out-of-sample
+ calibrated evaluation). A data card and a zipped bundle of the processed tables
ship with the repository, and every derived table carries lineage and
\texttt{rights\_reference} fields.

\section{Limitations and Ethics}
The system intentionally cannot ingest scraped or unlicensed media; researchers
remain responsible for confirming their own rights to any local footage. The
modeling results are illustrative and based on a small, highlight-only corpus;
they are not a predictive claim about match outcomes. No external execution or
real-money actions are performed.

\section{Conclusion and Future Work}
We presented a rights-safe, local-first soccer analytics pipeline that turns
permitted footage and open event data into detection datasets and match-outcome
models, with a hard, tested rights gate at every footage boundary. Future work
includes calibrating the detector on full-match footage and incorporating open
event features (e.g., xT, pressure) as stronger supervisory signal; the feature
scaler and calibrated classifier are already in place.

\bibliographystyle{plain}
\begin{thebibliography}{9}
\bibitem[StatsBomb]{statsbomb} StatsBomb. \emph{StatsBomb Open Data}. \url{https://github.com/statsbomb/open-data}.
\bibitem[Metrica]{metrica} Metrica Sports. \emph{Sample Game Data}. \url{https://github.com/metrica-sports/sample-data}.
\bibitem[SoccerNet]{soccernet} A. Cioppa et al. \emph{SoccerNet: A Scalable Dataset for Action Spotting in Soccer Videos}. CVPR 2020.
\bibitem[Redmon]{redmon2016yolo} J. Redmon, S. Divvala, R. Girshick, A. Farhadi. \emph{You Only Look Once: Unified, Real-Time Object Detection}. CVPR 2016.
\bibitem[Ultralytics]{ultralytics} Ultralytics. \emph{YOLOv8}. \url{https://github.com/ultralytics/ultralytics}.
\bibitem[scikit]{scikit} F. Pedregosa et al. \emph{Scikit-learn: Machine Learning in Python}. JMLR 2011.
\bibitem[PyTorch]{pytorch} A. Paszke et al. \emph{PyTorch: An Imperative Style, High-Performance Deep Learning Library}. NeurIPS 2019.
\end{thebibliography}

\end{document}
